\documentclass[11pt,a4paper]{article}
\usepackage[utf8]{inputenc}
\usepackage[vietnamese]{babel}
\usepackage{amsmath,amssymb}
\usepackage[version=4]{mhchem}
\usepackage{geometry}
\geometry{top=1.5cm,bottom=1.2cm,left=1.4cm,right=1.4cm,headheight=28pt,headsep=12pt,footskip=18pt}
\usepackage{graphicx}
\usepackage{xcolor}
\usepackage{tabularx}
\usepackage{array}
\usepackage{tikz}
\usepackage{fancyhdr}
\usepackage{enumitem}

\definecolor{hfRed}{HTML}{D32F2F}
\definecolor{hfGreen}{HTML}{8BC390}
\definecolor{mainGreen}{HTML}{2E7D32}
\definecolor{solBlue}{HTML}{1565C0}
\definecolor{tableGreen}{HTML}{E8F5E9}

\usepackage{ex_test}

\pagestyle{fancy}
\fancyhf{}
\renewcommand{\headrulewidth}{0.75pt}
\renewcommand{\headrule}{\hbox to\headwidth{\color{hfGreen}\leaders\hrule height \headrulewidth\hfill}}
\renewcommand{\footrulewidth}{0.75pt}
\renewcommand{\footrule}{\hbox to\headwidth{\color{hfGreen}\leaders\hrule height \footrulewidth\hfill}}

\fancyhead[L]{\small\itshape\color{hfRed}Tài liệu lưu hành nội bộ}
\fancyhead[R]{\small\color{hfRed}Thầy \textbf{TRẦN VĂN THẠNH} - 0777.470.803}

\fancyfoot[L]{\footnotesize\itshape\color{hfRed}CS1: 6/15 Nguyễn Hoàng, P. Kim Long, TP Huế \quad|\quad CS2: 24 Đặng Thái Thân, TP Huế}
\fancyfoot[R]{\small\textbf{\color{hfRed}Trang \thepage}}

\setlength{\parindent}{0pt}
\setlength{\parskip}{2pt}

\begin{document}

\begin{minipage}[t]{0.48\textwidth}
\textbf{SỞ GD\&ĐT THÀNH PHỐ HUẾ}\\
\textbf{TRƯỜNG THPT HAI BÀ TRƯNG}\\
\textit{ĐỀ CHÍNH THỨC}\\
\textit{(Đề thi có 03 trang)}
\end{minipage}
\hfill
\begin{minipage}[t]{0.48\textwidth}
\raggedleft
\textbf{\color{hfRed}KIỂM TRA GIỮA KÌ 1 NĂM HỌC 2026 - 2027}\\
\textbf{MÔN HÓA HỌC -- Khối lớp 10}\\
\textit{Thời gian làm bài: 45 phút (không kể thời gian giao đề)}\\
\textbf{MÃ ĐỀ: 132}
\end{minipage}

\vspace{4pt}
\noindent
\textit{Họ, tên thí sinh:} \dotfill \textit{Lớp:} \dotfill \textit{Số báo danh:} \dotfill

\vspace{3pt}
\noindent
\textit{\footnotesize Cho nguyên tử khối (amu): H = 1; C = 12; N = 14; O = 16; Si = 28; P = 31; S = 32; Li = 7; Na = 23; K = 39; Be = 9; Mg = 24; Ca = 40; Ba = 137.\\
Cho số hiệu nguyên tử: H: 1; N: 7; C: 6; Cl: 17; O: 8; S: 16; Li: 3; Na: 11; K: 19; Be: 4; Mg: 12; Ca: 20.\\
(Học sinh không sử dụng bảng hệ thống tuần hoàn các nguyên tố hóa học)}

\vspace{3pt}
\hrule color{hfGreen} height 0.5pt

\vspace{4pt}
\noindent
{\color{mainGreen}\textbf{\Large A. PHẦN TRẮC NGHIỆM: (7 điểm)}}

\vspace{3pt}
\noindent
{\color{mainGreen}\textbf{\large PHẦN I (3 điểm). TRẮC NGHIỆM NHIỀU LỰA CHỌN}}\\
\textit{\small Thí sinh trả lời từ câu 1 đến câu 12. Mỗi câu hỏi thí sinh chỉ chọn một phương án.}
\vspace{2pt}

\begin{ex}
Silicon được dùng trong công nghệ sản xuất chip máy tính hiện đại.. Vị trí của nguyên tố silicon trong bảng tuần hoàn thuộc ô 14, chu kì 3, nhóm IVA. Hãy chọn nhận đúng về silicon.
\choice
{\True Cấu hình electron nguyên tử của nguyên tố silicon là $\ce{[Ne] 3s^2 3p^2}$}
{Lớp N có 4 electron hoá trị}
{Số điện tích hạt nhân của silicon bằng 28}
{Silic nằm ở hàng thứ 4}
\end{ex}
\begin{ex}
Hầu hết các hạt nhân nguyên tử được tạo nên từ các hạt
\choice
{electron, proton và neutron}
{electron và neutron}
{\True proton và neutron}
{electron và proton}
\end{ex}
\begin{ex}
Trong số những quá trình kể dưới đây, cho biết đâu là hiện tượng vật lí ?
\choice
{Dẫn khí carbon dioxide vào nước vôi trong, làm nước vôi trong vẩn đục}
{Cho vôi sống ($\ce{CaO}$) vào nước được $\ce{Ca(OH)2}$}
{\True Nước bị hoá hơi khi đun nóng}
{Đốt cháy đường mía tạo thành chất màu đen và có mùi khét}
\end{ex}
\begin{ex}
Số phân lớp electron có trong lớp O là
\choice
{2}
{4}
{3}
{\True 5}
\end{ex}
\begin{ex}
Phát biểu nào sau đây không đúng?
\choice
{Số khối của nguyên tử bằng tổng số proton và số neutron}
{Nguyên tố hoá học là những nguyên tử có cùng số đơn vị điện tích hạt nhân}
{Số hiệu nguyên tử bằng số đơn vị điện tích hạt nhân nguyên tử}
{\True Đồng vị là các nguyên tố có cùng số proton nhưng khác nhau số neutron}
\end{ex}
\begin{ex}
Phát biểu nào sau đây đúng?
\choice
{Electron ở orbital 3p có mức năng lượng thấp hơn electron ở orbital 3s}
{Những electron ở gần hạt nhân có mức năng lượng cao nhất}
{\True Những electron ở lớp K có mức năng lượng cao nhất}
{Các electron trong cùng một lớp có năng lượng bằng nhau}
\end{ex}
\begin{ex}
Tổng số các hạt proton, nơtron và electron của một nguyên tử nguyên tố X là 28. Trong đó số hạt mang điện nhiều hơn số hạt không mang điện là 8. X là:
\choice
{Cl}
{Al}
{\True F}
{Si}
\end{ex}
\begin{ex}
Đồng (Copper) là kim loại dễ dát mỏng và uốn lượn, có khả năng dẫn điện và dẫn nhiệt tốt... Trong tự nhiên, nguyên tố Copper có hai đồng vị là $^{63}\ce{Cu}$ và $^{65}\ce{Cu}$. Nguyên tử khối trung bình của Copper là 63,54. Thành phần \% số nguyên tử của đồng vị $^{63}\ce{Cu}$ là
\choice
{27\%}
{\True 73\%}
{50\%}
{54\%}
\end{ex}
\begin{ex}
Phát biểu nào sau đây về chu kì là đúng?
\choice
{Dãy các nguyên tố mà nguyên tử của chúng có cùng số lớp electron, được xếp theo chiều số neutron tăng dần}
{\True Mở đầu mỗi chu kì là kim loại điển hình và kết thúc là một khí hiếm (trừ chu kì 1)}
{Dãy các nguyên tố mà nguyên tử của chúng có cùng số lớp electron, được xếp theo chiều khối lượng nguyên tử tăng dần}
{Dãy các nguyên tố mà nguyên tử của chúng có cùng số lớp electron, được xếp theo chiều số khối tăng dần}
\end{ex}
\begin{ex}
X là nguyên tố hóa học có trong thành phần của chất có tác dụng oxi hóa và sát khuẩn cực mạnh, thường được sử dụng với mục đích khử trùng và tẩy trắng trong lĩnh vực thủy sản, dệt nhuộm, xử lí nước cấp, nước thải, nước bể bơi. Nguyên tử X có số hiệu nguyên tử $Z = 17$. Số electron lớp hoá trị trong nguyên tử X là
\choice
{\True 7}
{6}
{5}
{4}
\end{ex}
\begin{ex}
Nội dung nào dưới đây không phải là đối tượng nghiên cứu của hóa học?
\choice
{Ứng dụng của chất}
{Tính chất và sự biến đổi của chất}
{Thành phần cấu trúc của chất}
{\True Sự lớn lên và sinh sản của tế bào}
\end{ex}
\begin{ex}
Cho dãy chất được biểu diễn bằng công thức hóa học như sau: $\ce{HNO3, O3, NaOH, He, Cl2, CuO, H2SO4, NaBr, Zn, CH4, MnO2}$. Số hợp chất trong dãy trên là
\choice
{\True 7}
{5}
{4}
{6}
\end{ex}
\vspace{4pt}
\noindent
{\color{mainGreen}\textbf{\large PHẦN II (2 điểm). CÂU TRẮC NGHIỆM ĐÚNG / SAI}}\\
\textit{\small Thí sinh trả lời từ câu 1 đến câu 2. Trong mỗi ý a), b), c), d) ở mỗi câu, thí sinh chọn đúng (Đ) hoặc sai (S).}
\vspace{2pt}

\begin{ex}
Các phát biểu sau đây là đúng hay sai?
\choiceTF
{\True Proton là hạt mang điện tích dương, nằm trong hạt nhân, khối lượng gần bằng 1 amu}
{\True Nguyên tử trung hòa điện, có kích thước lớn hơn nhiều so với hạt nhân, nhưng có khối lượng gần bằng khối lượng hạt nhân}
{Nguyên tử R có điện tích lớp vỏ là $-40,6\cdot 10^{-19}\text{ C}$ thì hạt nhân nguyên tử R có 26 proton}
{Nguyên tử X có 32 electron và 35 neutron, kí hiệu nguyên tử X là $^{16}_{32}\ce{S}$}
\end{ex}
\begin{ex}
Nguyên tử iron (Fe) có kí hiệu là $^{56}_{26}\ce{Fe}$. Cho các phát biểu sau:
\choiceTF
{\True Nguyên tử Fe có 2 electron lớp ngoài cùng}
{\True Fe là nguyên tố d}
{Trong bảng tuần hoàn Fe thuộc nhóm IIA}
{Trong nguyên tử Fe số hạt không mang điện chiếm 39,6\%. Về khối lượng nguyên tử}
\end{ex}
\vspace{4pt}
\noindent
{\color{mainGreen}\textbf{\large PHẦN III (2 điểm). TRẮC NGHIỆM TRẢ LỜI NGẮN}}\\
\textit{\small Thí sinh trả lời từ câu 1 đến câu 4. Viết câu trả lời vào ô quy định.}
\vspace{2pt}

\begin{ex}
Cho các phát biểu sau:\\
(1) Tổng số proton và số neutron trong hạt nhân bằng giá trị số khối.\\
(2) Trong một nguyên tử luôn có số proton bằng số neutron.\\
(3) Các đồng vị có tính chất vật lí khác nhau nhưng tính chất hoá học giống nhau.\\
(4) Số proton bằng số đơn vị điện tích hạt nhân nguyên tử.\\
(5) Trong nguyên tử, hạt electron có khối lượng không đáng kể so với các hạt còn lại.\\
Trong các phát biểu trên có bao nhiêu phát biểu đúng?

\noindent\textit{Trả lời:} \dotfill
\end{ex}
\begin{ex}
Trong bảng tuần hoàn nguyên tố X thuộc chu kì 4, nhóm IIA. Số hiệu nguyên tử của X là bao nhiêu?

\noindent\textit{Trả lời:} \dotfill
\end{ex}
\begin{ex}
Các electron của nguyên tử nguyên tố X được phân bố trên ba lớp, lớp thứ ba có 5 electron. Số đơn vị điện tích hạt nhân của nguyên tử nguyên tố X là bao nhiêu?

\noindent\textit{Trả lời:} \dotfill
\end{ex}
\begin{ex}
Nước cất ($\ce{H2O}$) là nước tinh khiết, nguyên chất, được điều chế bằng cách chưng cất và thường được sử dụng trong y tế như pha chế thuốc tiêm, thuốc uống, biệt dược, rửa dụng cụ y tế, rửa vết thương,... Biết phân tử này chỉ được tạo nên từ các đồng vị $^1_1\ce{H}$ và $^{16}_8\ce{O}$. Tổng số hạt electron, proton và neutron trong một phân tử $\ce{H2O}$ là bao nhiêu?

\noindent\textit{Trả lời:} \dotfill
\end{ex}
\vspace{6pt}
\noindent
{\color{mainGreen}\textbf{\Large B. PHẦN TỰ LUẬN: (3 điểm)}}
\vspace{2pt}

\noindent
\textbf{Câu 1: (1,0 điểm)} Nguyên tử của nguyên tố X nằm ở ô thứ 17, chu kì 3 nhóm VIIA trong bảng HTTH các nguyên tố hoá học.\\
a) Viết cấu hình electron nguyên tử X.\\
b) Viết công thức oxide cao nhất đối với oxigene và hợp chất với Hidrogene.

\vspace{1pt}
\noindent
\multido{\n=1+1}{6}{\dotfill\\[4pt]}

\vspace{4pt}
\noindent
\textbf{Câu 2: (1,0 điểm)} Nguyên tử khối trung bình của nguyên tố R là 79,91. R có hai đồng vị. Biết $^{79}\ce{R}$ chiếm 54,5\%. Nguyên tử khối của đồng vị thứ 2 bằng bao nhiêu?

\vspace{1pt}
\noindent
\multido{\n=1+1}{6}{\dotfill\\[4pt]}

\vspace{4pt}
\noindent
\textbf{Câu 3: (1,0 điểm)} Một loại quặng chứa 2 muối cacbonate của 2 kim loại kế tiếp nhau trong nhóm IIA được dùng trong vật liệu xây dựng, luyện kim, cải tạo đất chua... Cho 10,80 g quặng trên tác dụng với dd $\ce{H2SO4}$ loãng dư. Chất khí thu được cho hấp thụ hoàn toàn vào dd $\ce{Ba(OH)2}$ dư thì thu được 23,64 g kết tủa.\\
a. Viết các phương trình phản ứng hoá học xảy ra. Tìm công thức 2 muối trên.\\
b. Cần dùng bao nhiêu ml dung dịch HCl 30\%, $d = 1,5\text{ gam/ml}$ để phản ứng hết 5,4 gam quặng trên.

\vspace{1pt}
\noindent
\multido{\n=1+1}{8}{\dotfill\\[4pt]}

\vspace{6pt}
\begin{center}
\textbf{\large ---------- HẾT ----------}
\end{center}

\end{document}
